\section{Method and Theory}\label{sec:method}

\subsection{Exact reduction and the zero lift}\label{sec:reduction}

The only nonlinear constraint is $q^\top q=1$. Translation is eliminated exactly, not relaxed, so the nonconvex problem is a quartic on $\mathbb{S}^3$. ``Zero lift'' refers to the optimisation variable: no matrix variable is introduced at inference. The $10$-dimensional objects $v_2$ and $x_R$ are \emph{analysis} devices; they are never optimised over in the layer.

\begin{lemma}[Exact translation elimination]\label{lem:translation}
For $G\succ0$, $\min_t J_\theta(q,t)=J_{\mathrm{red}}(q;\theta)$, attained at $t^\star(q)=-G^{-1}Dv_2(q)$, and
\begin{equation}
\int_{\mathbb{R}^3}e^{-J_\theta(q,t)/\tau}\,dt=\frac{(\pi\tau)^{3/2}}{\sqrt{\det G}}\,e^{-J_{\mathrm{red}}(q;\theta)/\tau}.
\label{eq:gaussint}
\end{equation}
\end{lemma}

\begin{proof}
Completing the square gives $J_\theta(q,t)=(t-t^\star)^\top G(t-t^\star)+J_{\mathrm{red}}(q;\theta)$ with $t^\star=-G^{-1}Dv_2(q)$; the Gaussian integral over $\mathbb{R}^3$ is standard.
\end{proof}

\subsection{The rotation lift is an orthogonal image of the Veronese map}\label{sec:lift}

\begin{theorem}\label{thm:lift}
There exists $U\in O(10)$ such that for all $q\in\mathbb{R}^4$,
\begin{equation}
x_R(q)=2Uv_2(q),\qquad \|x_R(q)\|^2=4\|v_2(q)\|^2=4\|q\|^4 .
\label{eq:bridge}
\end{equation}
\end{theorem}

\begin{proof}[Proof sketch]
Every entry of $x_R$ is a quadratic form in $q$, and $v_2$ is a basis of $\mathrm{Sym}^2(\mathbb{R}^4)$; hence $x_R=Bv_2$ for a unique $B$. Since $R(q)$ is orthogonal, $\|x_R\|^2=\|q\|^4+3\|q\|^4=4\|q\|^4=4\|v_2\|^2$; as $\{v_2(q)\}$ spans $\mathbb{R}^{10}$, $B^\top B=4I$. Set $U=B/2$. The block structure is explicit: after permuting rows and columns, $U=\mathrm{diag}(\tfrac12H_4,\tfrac1{\sqrt2}H_2,\tfrac1{\sqrt2}H_2,\tfrac1{\sqrt2}H_2)$ with $H_4$ the Sylvester Hadamard matrix.
\end{proof}

\textbf{The homogeneous lift.} $x_R$ must be the \emph{homogeneous} lift \eqref{eq:lift}, whose first entry is $\|q\|^2$. With a literal constant $1$ in that slot the identity holds only on $\mathbb{S}^3$ and $U$ is not orthogonal off the sphere; the homogeneous lift makes Theorem~\ref{thm:lift} an identity in $\mathbb{R}^4$ and validates the conic argument of Section~\ref{sec:sdp}.

\subsection{The 21 \texorpdfstring{$SO(3)$}{SO(3)} constraints span the degree-two Veronese ideal}\label{sec:ideal}

\begin{theorem}\label{thm:ideal}
Let $A_i\in\mathbb{S}^{10}$ be the matrices of the 21 standard $SO(3)$ constraints: 3 column norms, 3 column inner products, 3 row norms, 3 row inner products, and 9 homogeneous cross-product equations $c_a\times c_b-x_0c_c=0$. Then
\begin{equation}
\dim\operatorname{span}\{A_i\}=20,
\label{eq:span20}
\end{equation}
and this span is exactly the degree-two vanishing ideal $\mathcal{I}_2(\nu_2(\mathbb{RP}^3))$ after the change of coordinates $x=2Uv_2$.
\end{theorem}

\begin{proof}[Proof sketch]
Each $A_i$ vanishes on $x_R(\mathbb{S}^3)$; the pullback is a homogeneous quartic vanishing on the sphere, hence identically zero, so $A_i$ lies in the kernel of $\phi:\mathbb{S}^{10}\to\mathbb{R}[q]_4$, $Q\mapsto Q(x_R(q))$. The map $\phi$ is surjective (the entries of $R_{ij}$ span all quadratics, and products of quadratics span all quartics), so $\dim\ker\phi=55-35=20$. An exact rational rank computation gives $\operatorname{rank}\{A_i\}=20$; the only linear dependency is $\operatorname{tr}(R^\top R)=\operatorname{tr}(RR^\top)$. A 20-dimensional subspace of a 20-dimensional kernel is the kernel. The self-contained surjectivity and dimension argument, including the exact rank step, is proved in Appendix~\ref{app:dimcount}.
\end{proof}

\textbf{Dimension count and rank computation.} A dimension count alone, $55-\binom74=20$, is necessary but not sufficient, since it does not show that the 21 constraint matrices reach the count; the rank computation alone does not identify the span with the ideal. Together they do. The count is also sensitive to the constraint set: six off-diagonal inner products alone give rank 15, not 20, which is why the constraint deficit of the released 24-constraint implementation (Section~\ref{sec:thirteen}) is not a bookkeeping detail.

\subsection{The rotation SDP is a Veronese relaxation}\label{sec:sdp}

\begin{theorem}\label{thm:sdp}
Let $\{B_k\}_{k=1}^{20}$ be a basis of $\mathcal{I}_2$, and let $M\in\mathbb{S}^{10}$. The SDP
\begin{equation}
\min_{Y\succeq0}\ \langle M,Y\rangle\quad\text{s.t.}\quad \langle B_k,Y\rangle=0\ (k=1,\dots,20),\quad \operatorname{tr}Y=1
\label{eq:sdpY}
\end{equation}
is the Shor (Veronese) relaxation of $\min_{q\in\mathbb{S}^3}v_2(q)^\top Mv_2(q)$. In $x$-coordinates, with $Q=\tfrac14UMU^\top$, it is equivalent to
\begin{equation}
\min_{X\succeq0}\ \langle Q,X\rangle\quad\text{s.t.}\quad X_{00}=1,\quad \langle A_i,X\rangle=0,
\label{eq:sdpX}
\end{equation}
under $X=4UYU^\top$.
\end{theorem}

\begin{proof}[Proof sketch]
The two PSD cones are related by an orthogonal congruence. On $(\mathcal{I}_2)^\perp$, $\operatorname{tr}Y=\langle ee^\top,Y\rangle$ with $e=(1,1,1,1,0,\dots,0)^\top$, because both sides agree on the rank-one generators $v_2(q)v_2(q)^\top$; since $U^\top e_0=\tfrac12e$, $X_{00}=\langle ee^\top,Y\rangle$. The objectives agree by orthogonality. The dual is $\max\mu$ s.t. $M-\mu I-\sum_k\lambda_kB_k\succeq0$, an SOS-modulo-ideal certificate.
\end{proof}

\textbf{Scope of the identification.} This is the rotation subproblem after exact translation elimination. The SDP is the degree-two Veronese/Shor relaxation, with the single trace/$X_{00}$ condition as the localising constraint rather than a moment matrix of a polynomial ideal; higher-order relaxations are outside the analysis. The relation to the released $13\times13$ implementation is Theorem~\ref{thm:thirteen}.

\subsection{A constructed instance with a certified gap}\label{sec:example}

\begin{example}\label{ex:choilam}
Let
\begin{equation}
C(q)=q_0^2q_1^2+q_1^2q_2^2+q_2^2q_0^2+q_3^4-4q_0q_1q_2q_3\ge0,
\label{eq:choilam}
\end{equation}
the Choi--Lam nonnegative non-SOS quartic \cite{choi1977extremal,choi1980zeros}, and let $M=M_{CL}+2I_{10}$, where $M_{CL}$ is the symmetric matrix of $C$ in the $v_2$ basis. The true minimum is $f^\star=2$ (since $C\ge0$ and $\|v_2\|^2=1$). Solving the $10\times10$ rotation SDP \eqref{eq:sdpX} gives $f^\star_{\mathrm{SDP}}=1.9658119658$, with a dual certificate $Z=Q-\mu E_{00}-\sum_i\lambda_iA_i$ satisfying $\lambda_{\min}(Z)=-1.8\times10^{-11}$ and $\mu=1.9658119658$. The certified gap is therefore at least $0.034188$; the primal solution has $\operatorname{rank}(X^\star)=6$ and $\lambda_1/\lambda_2=1.451$.
\end{example}

This is our own construction; a different perturbation changes the numerical gap and eigen-ratios, so we report the certified instance above, and Table~\ref{tab:sdp_cert} lists every quantity. ``Certified'' here means that the dual matrix is PSD at solver tolerance, not a symbolic proof. Whether $M$ is realisable as the Gram matrix of an actual PnP or matrix-weighted correspondence set is an open question; the physical six-point instance of Section~\ref{sec:thirteen} is the complementary, realisable separation result. The existence of nonnegative non-SOS forms is classical \cite{hilbert1888,blekherman2006}, and the Choi--Lam form is the standard explicit witness; the point of Example~\ref{ex:choilam} is not that such forms exist, but that the degree-two Veronese relaxation provably fails on a cost in the QPEP class.

\subsection{The \texorpdfstring{$3\times3$}{3x3} Riemannian implicit differentiation}\label{sec:implicit}

\begin{theorem}\label{thm:implicit}
Let $J(q)=v_2(q)^\top Mv_2(q)$ and let $q^\star$ be a stationary point on $\mathbb{S}^3$. Then:
\begin{enumerate}
\item $\nabla J(q^\star)=\lambda^\star q^\star$ with $\lambda^\star=q^{\star\top}\nabla J(q^\star)=4J(q^\star)$.
\item The Riemannian Hessian in the basis $E$ is
\begin{equation}
H_{\mathbb{S}^3}(q^\star)=E(q^\star)^\top\big(\nabla^2J(q^\star)-4J(q^\star)I_4\big)E(q^\star),
\label{eq:hess}
\end{equation}
where $\nabla^2J=2J_v^\top MJ_v+2\sum_{a=1}^{10}(Mv_2)_aH_{v,a}$.
\item If $H_{\mathbb{S}^3}(q^\star)\succ0$, then the local minimiser is locally $C^\infty$ in $\theta$ and
\begin{align}
\frac{\partial q^\star}{\partial\theta}&=-E H_{\mathbb{S}^3}^{-1}E^\top\Big(2J_v^\top\frac{\partial M}{\partial\theta}v_2\Big),\label{eq:dq}\\
\frac{\partial t^\star}{\partial\theta}&=-G^{-1}\Big(\frac{\partial G}{\partial\theta}t^\star+\frac{\partial D}{\partial\theta}v_2+DJ_v\frac{\partial q^\star}{\partial\theta}\Big).\label{eq:dt}
\end{align}
\item The map is antipodally invariant: $H_{\mathbb{S}^3}(-q^\star)=H_{\mathbb{S}^3}(q^\star)$, $\partial(-q^\star)/\partial\theta=-\partial q^\star/\partial\theta$, and $\partial\,\mathrm{vec}(R(q^\star))/\partial\theta$ is identical for $\pm q^\star$.
\end{enumerate}
\end{theorem}

\begin{proof}[Proof sketch]
Euler homogeneity gives the multiplier. Differentiating $F=E^\top\nabla J=0$ along $\delta q=E\delta\xi$ and using $(\delta E)^\top q=-E^\top\delta q$ gives the Hessian \eqref{eq:hess}. The implicit Jacobians \eqref{eq:dq}--\eqref{eq:dt} follow by differentiating $F(q^\star(\theta),\theta)=0$ and $t^\star=-G^{-1}Dv_2$, i.e.\ from the implicit-function theorem for the reduced stationarity system \cite{dontchev2009}. Antipodal invariance follows from $v_2(-q)=v_2(q)$, $J_v(-q)=-J_v(q)$ and $E(-q)=-E(q)$.
\end{proof}

\textbf{Factor conventions.} The conventions are a common source of disagreement between implementations. $\delta q=E\delta\xi$ corresponds to a rotation vector $2\delta\xi$, so $H_\omega=\tfrac14H_{\mathbb{S}^3}$, and a covariance in one convention differs from the other by a factor of 16 in the eigenvalues. QPEP's covariance analysis \cite{wu2022qpep} uses its own parameterisation, and comparisons should be made up to this constant. The backward pass \eqref{eq:dq} is a single $3\times3$ positive-definite solve per active mode and touches the $10\times10$ data matrix only through the fixed $4\times10$ Jacobian $J_v$.

\subsection{A differentiability certificate}\label{sec:cert}

\begin{theorem}[Sufficient condition]\label{thm:cert}
Suppose the global minimiser is unique up to $q\leftrightarrow-q$, $G\succ0$, $H_{\mathbb{S}^3}(q^{(1)})\succ0$, and $\Delta J=J(q^{(2)})-J(q^{(1)})>0$ for the next-best distinct $SO(3)$ class. Then there is a neighbourhood of $\theta$ on which the global pose map is $C^\infty$ and its derivative is given by Theorem~\ref{thm:implicit}.
\end{theorem}

\begin{proof}
The implicit-function theorem gives a smooth local minimiser; continuity of $\Delta J$ keeps it globally optimal; $G\succ0$ makes $t^\star$ unique.
\end{proof}

\textbf{Two counterexamples delimit the theorem.}
(i) $\Delta J=0$ does not preclude a smooth global \emph{selection}: with $J_1\equiv0$ and $J_2(\theta)=(\theta-\theta_0)^2$, branch 1 is a smooth global minimiser for all $\theta$, although the two branches tie at $\theta_0$. The correct statement concerns uniqueness of the global minimiser, not differentiability of every selection.
(ii) $H_{\mathbb{S}^3}\succ0$ is not necessary: for $J(x;\theta)=x^4+\theta^2x^2$, the unique minimiser $x^\star(\theta)\equiv0$ is smooth at $\theta=0$ although the Hessian vanishes. The sharp necessary condition is that the reduced stationarity system admits a differentiable single-valued selection that is globally optimal at every point; $H\succ0$ is the standard sufficient condition \cite{dontchev2009}.

\textbf{A QPEP-specific counting issue.} A naive count gives $K\le27$ real minima, which has no cited source and is incompatible with the free antipodal symmetry $(q,\lambda)\mapsto(-q,\lambda)$ for a finite root set. The released LibQPEP solver returns up to 40 candidate complex roots, and a total-degree homotopy confirms 40 complex projective roots (80 affine roots with the $\pm$ double cover) on $4/4$ instances (Section~\ref{sec:rootcount}), with the real projective count varying by instance. We therefore write ``$K$ enumerated real minima, at most 40 candidates in the released solver'', treat enumeration completeness as assumption (A2), and verify residuals numerically.

\subsection{Multi-modal Laplace smoothing}\label{sec:laplace}

\begin{theorem}\label{thm:laplace}
Let $q^{(k)}$, $k=1,\dots,K$, be the enumerated nondegenerate local minima of $J_{\mathrm{red}}$ on $\mathbb{S}^3/\{\pm1\}$, with $H_k=H_{\mathbb{S}^3}(q^{(k)})\succ0$. For the Boltzmann posterior $p_\theta(q,t)\propto e^{-J_\theta(q,t)/\tau}$:
\begin{enumerate}
\item \textbf{Translation is exact.} $\int_{\mathbb{R}^3}e^{-J_\theta(q,t)/\tau}dt=(\pi\tau)^{3/2}(\det G)^{-1/2}e^{-J_{\mathrm{red}}(q;\theta)/\tau}$.
\item \textbf{Rotation is Laplace-approximated.} With the Haar measure $\int_{SO(3)}dR=8\pi^2$ \cite{chirikjian2012},
\begin{equation}
\begin{split}
&\int_{SO(3)}e^{-J_{\mathrm{red}}/\tau}dR\\
&\quad=8(2\pi\tau)^{3/2}\sum_{k=1}^K\frac{e^{-J_k/\tau}}{\sqrt{\det H_k}}\big(1+O(\tau)\big),
\end{split}
\label{eq:rotationint}
\end{equation}
and the mode weights are
\begin{equation}
w_k=\operatorname{softmax}_k\Big(-\frac{J_k}{\tau}-\frac12\log\det H_k\Big).
\label{eq:weights}
\end{equation}
\item \textbf{Gradients.} For a downstream loss $\ell$,
\begin{equation}
\begin{split}
\nabla_\theta\bar\ell={}&\sum_kw_k\nabla_\theta\ell_k\\
&+\operatorname{Cov}_w\Big(\ell_k,\ -\frac1\tau\nabla_\theta J_k-\frac12\nabla_\theta\log\det H_k\Big),
\end{split}
\label{eq:smgrad}
\end{equation}
with $\nabla_\theta J_k=v_k^\top(\partial_\theta M)v_k$ and
\begin{equation}
\begin{split}
\nabla_\theta\log\det H_k&=\operatorname{tr}\Big(H_k^{-1}\frac{dH_k}{d\theta}\Big),\\
\frac{dH_k}{d\theta}&=\frac{\partial H}{\partial\theta}\Big|_{q^{(k)}}+\frac{\partial H}{\partial q}\Big|_{q^{(k)}}\frac{\partial q^{(k)}}{\partial\theta}.
\end{split}
\label{eq:dlogdet}
\end{equation}
The second term contains third $q$-derivatives of $J$.
\item \textbf{Temperature limit.} As $\tau\to0$ with $\Delta J>0$, the posterior mean and the soft argmin converge to the deterministic global minimiser, and their gradients converge to Theorem~\ref{thm:implicit}. The KL loss itself, $\mathcal{L}_{\mathrm{KL}}=\tau^{-1}J_{\mathrm{gt}}+\log Z_{\mathrm{Laplace}}$, has a gradient that scales as $1/\tau$ and does not converge without a $\tau$-rescaling. The rescaled (annealed) loss $\tau\mathcal{L}_{\mathrm{KL}}=J_{\mathrm{gt}}+\tau\log Z_{\mathrm{Laplace}}$ is analysed in Theorem~\ref{thm:annealed}: its gradient is finite at every $\tau>0$ and has a curvature-weighted Clarke limit as $\tau\to0^+$.
\end{enumerate}
\end{theorem}

\begin{proof}[Proof sketch]
The translation integral is Lemma~\ref{lem:translation}. For rotation, expand $J$ to second order around each $q^{(k)}$ in the normal chart $q(u)=\mathrm{normalize}(q^{(k)}+E_ku)$; the Gaussian integral gives $(2\pi\tau)^{3/2}(\det H_k)^{-1/2}$. The local Jacobian is $dR=8\,dq$ per sheet; the two sheets of $\pm q$ and the double-cover factor $1/2$ combine to $8$ per $SO(3)$ mode. The standard Laplace expansion gives relative error $O(\tau)$ \cite{bleistein1986,wong1989}. The gradient formula is the chain rule with the envelope theorem for $J_k$ and the softmax Jacobian for $w_k$. Appendix~\ref{app:haar} derives the Jacobian $dR=8\,du$, the double-cover cancellation, the absolute constant of $Z_{\mathrm{Lap}}$, and the full $\theta$-derivative of $\log\det H_k$ used below.
\end{proof}

\textbf{Smoothness at branch crossings.} The \emph{exact} Boltzmann posterior is $C^\infty$ in $\theta$ for every $\tau>0$, because $\mathbb{S}^3$ is compact and $J_\theta$ is smooth (by dominated convergence, without requiring $\Delta J>0$). The \emph{Laplace approximation} is not uniformly accurate at a crossing: as $\Delta J\to0$ and an eigenvalue of $H_k$ tends to zero, the Gaussians overlap and the $1/\sqrt{\det H_k}$ weights diverge, so a uniform (Airy) expansion or a one-dimensional numerical integral in the merging direction is required \cite{bleistein1986}. The smoothed layer is therefore smooth for every $\tau>0$, while the Laplace implementation is $O(\tau)$-accurate away from bifurcations.

\textbf{Normalisation.} The constants matter for $Z$ but cancel in $w_k$. With the probability Haar measure $d\mu_{\mathrm{Haar}}=dR/(8\pi^2)$ on $SO(3)$, the absolute Laplace partition function is
\begin{equation}
Z_{\mathrm{Lap}}(\theta)=\frac{2\sqrt2\,\pi\tau^3}{\sqrt{\det G}}\sum_{k=1}^K\frac{e^{-J_k/\tau}}{\sqrt{\det H_k}},
\label{eq:zlapabs}
\end{equation}
whereas the unnormalised measure $dR$ multiplies \eqref{eq:zlapabs} by $8\pi^2$. Appendix~\ref{app:haar} derives both constants from $dR=8\,du$ and the double cover, and shows that the two conventions give the same $w_k$ and the same $\theta$-gradient. The common normalisation $Z_{\mathrm{Lap}}=(\pi\tau)^3(\det G)^{-1/2}\sum_ke^{-J_k/\tau}/\sqrt{\det H_k}$ is neither of these absolute constants, so it should not be quoted as $Z$; it differs from \eqref{eq:zlapabs} by a $\theta$-independent factor at fixed $\tau$. If rotation-vector coordinates $\omega=2\delta\xi$ are used, $\det H_\omega=\det H_{\mathbb{S}^3}/64$, and one convention must be used consistently in \eqref{eq:weights}--\eqref{eq:dlogdet}.

\begin{theorem}[Annealed Laplace NLL: exact gradient and zero-temperature limit]\label{thm:annealed}
Let $q^{(k)}(\theta)$, $k=1,\dots,K$, be the nondegenerate minima of Theorem~\ref{thm:laplace}, with $J_k=J_\theta(q^{(k)},t^\star(q^{(k)}))$, $H_k\succ0$ and $G\succ0$, and let
\begin{equation}
\begin{split}
Z_{\mathrm{Lap}}(\theta)={}&C(\tau)\,(\det G(\theta))^{-1/2}\\
&\times\sum_{k=1}^K(\det H_k(\theta))^{-1/2}e^{-J_k(\theta)/\tau}
\end{split}
\label{eq:zlap}
\end{equation}
with $C(\tau)$ independent of $\theta$, and $w_k=(\det H_k)^{-1/2}e^{-J_k/\tau}/\sum_j(\det H_j)^{-1/2}e^{-J_j/\tau}$. Define the annealed Laplace NLL
\begin{equation}
L_{\tau\text{-KL}}(\theta):=J_\theta(q_{\mathrm{gt}},t_{\mathrm{gt}})+\tau\log Z_{\mathrm{Lap}}(\theta).
\label{eq:annealed}
\end{equation}
Then:
\begin{enumerate}
\item \textbf{(finite smooth gradient for every $\tau>0$)} including at exact branch ties $J_{k_1}=J_{k_2}$, $L_{\tau\text{-KL}}$ is locally $C^\infty$ in $\theta$ and
\begin{equation}
\begin{split}
\nabla_\theta L_{\tau\text{-KL}}={}&\nabla_\theta J_\theta(q_{\mathrm{gt}},t_{\mathrm{gt}})\\
&-\sum_{k=1}^Kw_k\Big(\nabla_\theta J_k+\frac{\tau}{2}\nabla_\theta\log\det H_k\Big)\\
&-\frac{\tau}{2}\nabla_\theta\log\det G .
\end{split}
\label{eq:annealedgrad}
\end{equation}
The outer factor $\tau$ cancels the $1/\tau$ in $\nabla_\theta z_k$ exactly, so the gradient stays finite as $\tau\to0$ instead of diverging like the unannealed KL loss.
\item \textbf{(zero-temperature limit)} with $\mathcal{K}_{\min}=\arg\min_kJ_k$ and curvature-volume weights $\bar w_k:=(\det H_k)^{-1/2}/\sum_{j\in\mathcal{K}_{\min}}(\det H_j)^{-1/2}$,
\begin{equation}
\lim_{\tau\to0^+}\nabla_\theta L_{\tau\text{-KL}}
=\nabla_\theta J_\theta(q_{\mathrm{gt}},t_{\mathrm{gt}})
-\sum_{k\in\mathcal{K}_{\min}}\bar w_k\,\nabla_\theta J_k,
\label{eq:clarke}
\end{equation}
which is the curvature-weighted Clarke subgradient of $J_\theta(q_{\mathrm{gt}},t_{\mathrm{gt}})-\min_kJ_k(\theta)$. For a unique minimiser $k^\star$ it reduces to $\nabla_\theta J_{\mathrm{gt}}-\nabla_\theta J_{k^\star}$.
\end{enumerate}
\end{theorem}

\begin{proof}
Write $\tau\log Z_{\mathrm{Lap}}=\tau\log C(\tau)-\frac{\tau}{2}\log\det G+\tau\log\sum_kz_k$ with $z_k=(\det H_k)^{-1/2}e^{-J_k/\tau}$. Since $H_k\succ0$, the implicit-function theorem makes each branch and $H_k$ smooth in $\theta$, and $\sum_kz_k>0$ even when two $J_k$ coincide because the branches themselves stay nondegenerate. Using $\nabla_\theta z_k=z_k(-\tau^{-1}\nabla_\theta J_k-\frac12\nabla_\theta\log\det H_k)$ and $w_k=z_k/\sum_jz_j$ gives \eqref{eq:annealedgrad}; the $\tau$ in front of $\log Z$ cancels $\tau^{-1}$ exactly. For \eqref{eq:clarke}, factor $e^{-J_{\min}/\tau}$ out of $z_k$: $w_k\to(\det H_k)^{-1/2}/\sum_{j\in\mathcal{K}_{\min}}(\det H_j)^{-1/2}$ on $\mathcal{K}_{\min}$ and $w_k\to0$ elsewhere, because every $j\notin\mathcal{K}_{\min}$ has a strictly positive gap $\Delta_j=J_j-J_{\min}>0$. The terms $\frac{\tau}{2}\nabla_\theta\log\det H_k$ and $\frac{\tau}{2}\nabla_\theta\log\det G$ vanish, while $\nabla_\theta J_k$ is $\tau$-independent and bounded. The limit is a convex combination of the gradients of the active branches, hence lies in the Clarke subdifferential of the minimum. The envelope identity $\nabla_\theta J_k=v_2(q^{(k)})^\top(\partial_\theta M)v_2(q^{(k)})$ and the full derivative of $\log\det H_k$ are derived in Appendix~\ref{app:haar}.
\end{proof}

\subsection{Relation to the released \texorpdfstring{$13\times13$}{13x13} relaxation}\label{sec:thirteen}

The relation between the $10\times10$ rotation SDP and the released $13\times13$ SDP is not one equality or one inequality. It has three layers: an exact equivalence, an inequality hierarchy, and a strict physical separation.

\begin{theorem}[Three-part $10\times10$/$13\times13$ relation]\label{thm:thirteen}
Let $x_{13}=[x_R;t]$ with $t\in\mathbb{R}^3$ unconstrained,
$Q_{13}=\begin{bmatrix}Q_{RR}&Q_{Rt}\\Q_{Rt}^\top&G\end{bmatrix}$ with $G\succ0$, and $Q_{10}=Q_{RR}-Q_{Rt}G^{-1}Q_{Rt}^\top$. Write $f^\star_{10,\mathrm{full}}$, $f^\star_{13,\mathrm{full}}$ and $f^\star_{13,\mathrm{up}}$ for (a) the $10\times10$ value under the full 20-dimensional ideal, (b) the $13\times13$ value under the same full ideal, i.e.\ the upstream 24 constraints plus the three row-orthogonality equations (27 constraints, 20 dimensions), and (c) the $13\times13$ value under the released 24 constraints, whose span is 17-dimensional. Then:
\begin{enumerate}
\item \textbf{Full 20D equivalence (proved).} The degree-two vanishing ideal of $SE(3)$ in $\mathbb{S}^{13}$ is $\mathcal{I}_2(SE(3))=\mathcal{I}_2(SO(3))\oplus\mathbf{0}_{3\times3}$: dimension 20, with no coupling to $t$. Hence the block Schur-complement bijection $X^\star_{13,\mathrm{full}}=LX^\star_{10,\mathrm{full}}L^\top$, with $L=[I_{10};\,-G^{-1}Q_{Rt}^\top]$ of full column rank, gives
\begin{equation}
f^\star_{13,\mathrm{full}}\equiv f^\star_{10,\mathrm{full}},\quad\operatorname{rank}\big(X^\star_{13,\mathrm{full}}\big)\equiv\operatorname{rank}\big(X^\star_{10,\mathrm{full}}\big).
\label{eq:schur}
\end{equation}
This is an algebraic identity, not a numerical coincidence.
\item \textbf{The upstream 17D implementation is looser, strictly so in general (proved).} The released 24 constraints span only 17 of the 20 dimensions, so $\mathcal{V}_{17}\subsetneq\mathcal{I}_2(SO(3))$ and the upstream feasible set strictly contains the full-20D one. Therefore, for \emph{every} instance,
\begin{equation}
f^\star_{13,\mathrm{up}}\;\le\;f^\star_{13,\mathrm{full}}=f^\star_{10,\mathrm{full}}\;\le\;f^\star_{\mathrm{QPEP}},
\label{eq:hierarchy}
\end{equation}
the last inequality being the relaxation bound of Theorem~\ref{thm:sdp}. Because the argument is instance-independent, every full-20D instance that is not tight is automatically a non-tight instance of the released implementation.
\item \textbf{Physical 6-point separation (proved and independently reproduced).} Layer (2) is strict on a genuine matrix-weighted 3D--3D cost with positive weights and a unique ground-truth rotation, even where the full-20D SDP is exactly tight at rank one (Table~\ref{tab:6point}). By continuity the strict inequality persists on an open neighbourhood of that instance, so the separation is not a measure-zero pathology.
\end{enumerate}
\end{theorem}

\begin{proof}[Proof sketch]
(1) The degree-two vanishing ideal of $SE(3)$ in $\mathbb{S}^{13}$ is the ideal of $SO(3)$ in the $x_R$ block with the $t$ block free, so there is no coupling to $t$; the coupling in the cost is quadratic and is completed by the Schur complement, which is exactly the map $L$. Substituting $X_{13}=LX_{10}L^\top$ leaves the constraints and the objective invariant, and $L$ has full column rank because $G\succ0$. (2) A subspace inclusion is strict when the spanning set misses three dimensions; feasibility is monotone under a larger feasible set, and the last inequality in \eqref{eq:hierarchy} is the standard relaxation bound. (3) The instance is analysed in Section~\ref{sec:results_sdp}; its exact solution has $r_1^\top r_2=0$ at the full-ideal optimum and $r_1^\top r_2>0$ at the upstream optimum, so the inclusion is strict there, and strictness is an open condition.
\end{proof}

\textbf{Exact constraint comparison.} The 21 $SO(3)$ constraints of the $10\times10$ SDP span 20 dimensions (Theorem~\ref{thm:ideal}). The 24 upstream constraints of the $13\times13$ SDP span only 17, and projecting the 21 constraints one at a time onto that 17-dimensional span leaves exactly three nonzero residuals---the row orthogonality equations $r_1\cdot r_2=r_1\cdot r_3=r_2\cdot r_3=0$ (residuals 1.63--1.73; the other 18 are $\le6.2\times10^{-16}$). Row norms $\|r_i\|^2=h^2$ are already implied by the nine norm constraints together with column orthogonality, so the row dot products are the whole deficit. The $13\times13$ constraint generators were replicated row-by-row from the released SDPRLayers implementation and compared element-wise against the upstream generator bodies: all 24 constraints match (24/24) with max $|\Delta|=0$, and $\max_i|x^\top A_i x|=1.4\times10^{-16}$ on random rotations. This replication checks our reimplementation against the released generator bodies. The measured consequences are in Tables~\ref{tab:13v10} and~\ref{tab:6point}.

\subsection{Algorithm and complexity}\label{sec:algorithm}

\begin{algorithm}[t]
\caption{Forward and backward pass of \textsc{Laplace-QPEPLayer} (one batch element).}
\label{alg:layer}
\begin{algorithmic}[1]
\State \textbf{Forward:}
\State build $M_0,D,G$ from $\theta$;\ \ $M\gets M_0-D^\top G^{-1}D$
\State enumerate all real KKT roots $q^{(k)}$ of $\nabla J=4Jq$, $q^\top q=1$
\State keep local minima with $H_k\succ0$; sort by $J_k$; compute $w_k$ by \eqref{eq:weights}
\State \Return $(q^{(1)},t^\star(q^{(1)}))$ and/or $\{(q^{(k)},w_k)\}_{k=1}^{K}$
\State \textbf{Backward:}
\For{each active mode $k$}
  \State form $H_k$ by \eqref{eq:hess} and $dH_k/d\theta$ by \eqref{eq:dlogdet}
  \State solve $H_k\,x=E^\top\big(2J_v^\top(\partial_\theta M)v_2\big)$ and set $\partial q^{(k)}/\partial\theta=-Ex$
\EndFor
\State assemble the deterministic Jacobian \eqref{eq:dq}--\eqref{eq:dt} or the smoothed gradient \eqref{eq:smgrad}
\end{algorithmic}
\end{algorithm}

The nonconvex problem has 4 variables and 1 constraint. The dense linear algebra per mode is a $4\times10$ Jacobian, a $4\times4$ Euclidean Hessian and a $3\times3$ solve. The data-dependent cost is building $M$ (a sum over correspondences) and enumerating roots; both are independent of the SDP dimension, and neither scales with the number of modes. In the deterministic regime the backward pass is $O(1)$ per sample. In the smoothed regime it is $O(K)$ with $K$ the number of active modes, plus the third-derivative contraction for $\nabla\log\det H_k$ in \eqref{eq:dlogdet}. Root enumeration, not the $3\times3$ solves, dominates wall-clock time (Section~\ref{sec:compute}).
